\mychapter{Introduction}

\vspace{0.5cm}

% \subsection{Introduction}

 \subsection{Introduction}

The rapid advancement of information technology has transformed many sectors by replacing traditional paper-based workflows with efficient digital solutions. Universities have adopted digital systems for admission, registration, examination management, and result publication. However, the evaluation of descriptive admission examination scripts remains largely manual in many institutions. Since handwritten answer scripts require physical distribution, evaluation, moderation, and result compilation, the process becomes increasingly time-consuming and difficult to manage as the number of applicants continues to grow.

Public universities in Bangladesh receive thousands of admission applications every year for a limited number of seats. Descriptive examinations are typically evaluated by teachers from different disciplines such as Physics, Chemistry, Mathematics, and English. The conventional workflow requires physical distribution of answer scripts, manual allocation of questions, recording marks on paper, resolving discrepancies through moderation, and preparing results manually. Although this approach has been practiced for many years, it introduces several operational challenges when handling a large volume of answer scripts within a limited timeframe.

The manual evaluation process suffers from several limitations. Physical handling of scripts increases the risk of document misplacement, administrative errors, and processing delays. Since teachers often evaluate scripts containing student identification information, there is a possibility of conscious or unconscious bias during assessment. Moreover, maintaining consistent grading standards among multiple examiners is difficult, and administrators have limited capability to monitor evaluation progress or verify modifications made during the marking process.

Question-wise evaluation presents another significant challenge. Different subjects require specialized examiners, yet manually separating answer scripts according to individual questions and assigning them to appropriate teachers is repetitive and error-prone. Furthermore, comparing evaluations from multiple teachers and determining the final mark requires additional administrative effort, reducing transparency and efficiency throughout the moderation process.

Recent advances in document image processing, Optical Character Recognition (OCR), cloud computing, and modern web technologies provide opportunities to automate these tasks. \ac{OCR} enables extraction of printed text from scanned documents, while image processing techniques support segmentation of answer scripts into individual question regions. Combined with web-based evaluation platforms, these technologies enable remote evaluation and real-time monitoring of the assessment process.

Motivated by these challenges, this thesis presents the \textbf{Digital Admission Script Evaluation and Moderation System (DASEMS)}, a web-based platform that digitizes and automates the complete evaluation workflow for descriptive admission examinations. Unlike conventional systems that primarily focus on result management, \textbf{DASEMS} supports the entire evaluation lifecycle, including student information management, PDF upload, automated question segmentation, teacher assignment, answer evaluation, moderation, and final result generation.

The proposed system follows a role-based architecture consisting of \textbf{Super Administrator}, \textbf{Head Examiner}, and \textbf{Teacher} roles. To ensure fairness, student identities are concealed from evaluators using anonymous student codes. A Python-based document processing service automatically analyzes uploaded answer scripts, detects question boundaries, and extracts answer regions using a two-stage approach that combines text-layer analysis with \ac{OCR}-based fallback processing. Each answer is then assigned to multiple subject experts for independent evaluation. If significant differences occur among the assigned evaluators, the answer is automatically forwarded to the \textbf{Head Examiner} for moderation before finalizing the marks.

By integrating automated document processing, anonymous evaluation, role-based access control, and centralized moderation, \textbf{DASEMS} provides a comprehensive digital framework that improves efficiency, fairness, transparency, and accountability in the admission script evaluation process. Consequently, the proposed system offers a practical and scalable alternative to conventional paper-based evaluation for universities and other educational institutions conducting descriptive examinations.



\subsection{Background and Motivation}

University admission examinations play a crucial role in selecting qualified candidates for higher education. As the number of applicants continues to increase, universities must evaluate thousands of descriptive answer scripts within a limited time. Ensuring an efficient, fair, and transparent evaluation process has therefore become an important challenge for educational institutions.

Traditionally, admission answer scripts are physically distributed among teachers, evaluated on paper, moderated through meetings, and finally compiled into result sheets. Although this workflow has been used for many years, it requires considerable administrative effort and becomes increasingly inefficient when handling large volumes of answer scripts. Manual coordination, mark calculation, and moderation not only consume time but also increase the possibility of operational errors.

Maintaining consistency among multiple examiners is another significant challenge. Teachers may interpret marking criteria differently, resulting in variations in awarded marks. Moreover, because student identification is often visible during evaluation, complete anonymity cannot always be ensured, introducing the possibility of conscious or unconscious bias.

Advances in web technologies, cloud computing, document image processing, and Optical Character Recognition (OCR) provide opportunities to modernize this workflow. Digital platforms enable centralized management of evaluation activities, while OCR-assisted document analysis can automatically identify question boundaries and segment answer scripts into question-wise regions for evaluation. These technologies are intended to support human examiners by automating repetitive administrative tasks rather than replacing their academic judgment.

Motivated by these challenges, this research proposes the \textbf{Digital Admission Script Evaluation and Moderation System (DASEMS)}, which integrates secure web technologies, automated PDF processing, OCR-assisted question segmentation, anonymous evaluation, role-based access control, and centralized moderation into a unified platform. The proposed system aims to reduce administrative workload, improve examiner coordination, enhance fairness and transparency, and provide a scalable digital framework for descriptive admission script evaluation.


\subsection{Research Objectives}

The objective of this research is to develop a secure and efficient web-based platform for descriptive university admission script evaluation. The proposed system aims to automate answer script processing, anonymous evaluation, moderation, and result generation while improving the transparency, accuracy, and efficiency of the evaluation process.

\subsubsection{General Objective}

The general objective of this research is to design and develop a web-based Digital Admission Script Evaluation and Moderation System (DASEMS) for secure script processing, anonymous evaluation, moderation, and centralized result management.

\subsubsection{Specific Objectives}

The specific objectives of this research are as follows:

\begin{itemize}

\item To develop a role-based system for Super Administrator, Head Examiner, and Teacher.

\item To implement secure authentication and anonymous script evaluation.

\item To design a digital question bank for admission examinations.

\item To develop an OCR-based PDF processing module for question-wise answer extraction.

\item To implement automated teacher assignment and browser-based digital evaluation.

\item To develop a moderation mechanism for resolving significant differences in marks.

\item To generate final results and evaluation reports through a centralized management system.

\item To improve the efficiency, transparency, and fairness of the admission evaluation process.

\end{itemize}

\subsection{Scope of the Research}

This research focuses on the development of a web-based platform for managing descriptive university admission script evaluation. The proposed system automates the workflow from examination configuration and answer script processing to moderation and result generation while preserving human judgment during evaluation.

The platform supports centralized management of admission sessions, subjects, teachers, students, and evaluations. Administrators manage examinations and monitor progress, teachers evaluate assigned answers, and Head Examiners perform moderation when required.

A key contribution of the system is the automated processing of scanned answer scripts. Using a hybrid approach that combines embedded PDF text extraction and Optical Character Recognition (OCR), the system detects question boundaries and extracts question-wise answer regions for evaluation.

To ensure fairness, student identities are replaced with anonymous identifiers, and answers are automatically assigned to subject-specific teachers. A moderation mechanism forwards evaluations with significant mark differences to the Head Examiner for final review.

The scope of this research does not include automatic grading of handwritten answers. OCR is used only for document processing, while all marking decisions remain under the control of human examiners. The system has been evaluated using representative admission examination templates and provides a secure, scalable, and transparent framework for descriptive admission script evaluation.


\subsection{Research Contributions}

The major contributions of this research are summarized as follows:

\begin{itemize}

\item Design and development of a web-based platform for managing the complete workflow of descriptive university admission script evaluation and moderation.

\item Development of a secure role-based system supporting \textbf{Super Administrator}, \textbf{Head Examiner}, and \textbf{Teacher} roles, together with anonymous student identification to ensure fair evaluation.

\item Implementation of a two-stage PDF processing pipeline that combines embedded text extraction with \textbf{Optical Character Recognition (OCR)} for automatic question segmentation and answer extraction.

\item Automated distribution of question-wise answers to subject-specific teachers, along with a browser-based interface for digital annotation and mark submission.

\item Development of a moderation mechanism that detects significant differences between independent evaluations and forwards disputed answers to the \textbf{Head Examiner} for final review.

\item Design of a centralized, modular, and extensible architecture for managing evaluation records, progress monitoring, audit logs, and result generation, enabling efficient and transparent examination management.

\end{itemize}

\subsection{Organization of the Thesis}

The remainder of this thesis is organized into six chapters.

\textbf{Chapter 2} reviews the literature on digital examination systems, Optical Character Recognition (OCR), document image processing, online evaluation platforms, and role-based access control, highlighting the research gaps addressed in this work.

\textbf{Chapter 3} presents the proposed methodology, including the system architecture, database design, document processing pipeline, OCR-assisted question segmentation, teacher assignment, and moderation workflow.

\textbf{Chapter 4} describes the implementation of the proposed system, covering the software architecture, PDF processing module, evaluation interface, moderation module, and result generation.

\textbf{Chapter 5} presents the experimental results, system validation, performance analysis, and discussion of the effectiveness of the proposed system.

\textbf{Chapter 6} concludes the thesis by summarizing the research contributions, discussing current limitations, and suggesting future research directions.

Overall, the thesis progresses from the research background and literature review to the proposed methodology, implementation, experimental evaluation, and conclusion.




